% GENERATED FILE. Edit canonical lessons, then run npm run generate:books.
% canonical-source-hash: fnv1a64:9bbbc0c30db067e0
% canonical-lessons: ML-C466-syuttkesu, ML-C466-cumatu, ML-C466-porttar, ML-C466-klokkrum, ML-C466-valiccukontupokuka

\chapter{Bags and Loads}
\label{ch:ml-bags-and-loads}
\hlchaptermodality{466}

\begin{chapteropening}
\textbf{By the end of this chapter:} \emph{I can say a suitcase, a load, a porter, a cloakroom and to drag along.}

Complete the last lesson of chapter 466: I can say a suitcase, a load, a porter, a cloakroom and to drag along.
\end{chapteropening}

\section[\texorpdfstring{\ml{സ്യൂട്ട്കേസ്} (syūṭṭkēsŭ)}{സ്യൂട്ട്കേസ് (syūṭṭkēsŭ)}]{\ml{സ്യൂട്ട്കേസ്} (syūṭṭkēsŭ) — a suitcase}

\label{lesson:ML-C466-syuttkesu}



\begin{quote}
Before the new one: say the Malayalam for to wait standing, then the Malayalam for to run and board.
\end{quote}

\subsection*{You'll want to know: \ml{സ്യൂട്ട്കേസ്}}
\textbf{\ml{സ്യൂട്ട്കേസ്}} — \emph{syūṭṭkēsŭ} — \textquotedblleft{}a suitcase\textquotedblright{}.

\begin{quote}
\ml{എന്റെ} \ml{സ്യൂട്ട്കേസിന്} \ml{നല്ല} \ml{ഭാരമുണ്ട്}.
\end{quote}

\begin{quote}
\emph{enṟe syūṭṭkēsinŭ nalla bhāramuṇṭŭ.}
\end{quote}

\begin{quote}
My suitcase is very heavy.
\end{quote}

Weight is spoken of as something a thing has: \ml{ഭാരമുണ്ട്}, there is weight, it is heavy.

\begin{cousinweb}
\textbf{\ml{സ്യൂട്ട്കേസ്}} (\emph{syūṭṭkēsŭ}) was borrowed from English \emph{suitcase}, a case for carrying clothes on a journey.
\end{cousinweb}

\begin{grammarlens}[title={What you've built}]
One more word for getting around.
\end{grammarlens}

\subsection*{Your turn}
\begin{itemize}
\raggedright
  \item \emph{Say it:} \emph{syūṭṭkēsŭ}
  \item \emph{Say it:} \emph{syūṭṭkēsŭ}, once more
  \item \emph{Say it:} \emph{enṟe syūṭṭkēsinŭ nalla bhāramuṇṭŭ.} — My suitcase is very heavy.
  \item \emph{Recall:} say the Malayalam for to wait standing, then the Malayalam for to run and board, then say \emph{syūṭṭkēsŭ} again
\end{itemize}

\begin{tcolorbox}[breakable,colback=teal!4,colframe=teal!35!black,title={Before you move on}]
What does \ml{സ്യൂട്ട്കേസ്} mean? (\textquotedblleft{}A suitcase\textquotedblright{}.) Say it once more.
\end{tcolorbox}
\section[\texorpdfstring{\ml{ചുമട്} (cumaṭŭ)}{ചുമട് (cumaṭŭ)}]{\ml{ചുമട്} (cumaṭŭ) — a load; a burden}

\label{lesson:ML-C466-cumatu}



\begin{quote}
Before the new one: say the Malayalam for to run and board, then the Malayalam for a suitcase.
\end{quote}

\subsection*{You'll want to know: \ml{ചുമട്}}
\textbf{\ml{ചുമട്}} — \emph{cumaṭŭ} — \textquotedblleft{}a load; a burden\textquotedblright{}.

\begin{quote}
\ml{അയാൾ} \ml{തലയിൽ} \ml{വലിയ} \ml{ചുമട്} \ml{ചുമന്നു}.
\end{quote}

\begin{quote}
\emph{ayāḷ talayil valiya cumaṭŭ cumannu.}
\end{quote}

\begin{quote}
He carried a big load on his head.
\end{quote}

The noun of \ml{ചുമക്കുക}, to carry: what you carry is a \ml{ചുമട്}.

\begin{cousinweb}
\textbf{\ml{ചുമട്}} (\emph{cumaṭŭ}) is native Dravidian, from the root of \ml{ചുമക്കുക}, to carry a burden (DEDR 2677).
\end{cousinweb}

\begin{grammarlens}[title={What you've built}]
One more word for getting around.
\end{grammarlens}

\subsection*{Your turn}
\begin{itemize}
\raggedright
  \item \emph{Say it:} \emph{cumaṭŭ}
  \item \emph{Say it:} \emph{cumaṭŭ}, once more
  \item \emph{Say it:} \emph{ayāḷ talayil valiya cumaṭŭ cumannu.} — He carried a big load on his head.
  \item \emph{Recall:} say the Malayalam for to run and board, then the Malayalam for a suitcase, then say \emph{cumaṭŭ} again
\end{itemize}

\begin{tcolorbox}[breakable,colback=teal!4,colframe=teal!35!black,title={Before you move on}]
What does \ml{ചുമട്} mean? (\textquotedblleft{}A load; a burden\textquotedblright{}.) Say it once more.
\end{tcolorbox}
\section[\texorpdfstring{\ml{പോർട്ടർ} (pōrṭṭar)}{പോർട്ടർ (pōrṭṭar)}]{\ml{പോർട്ടർ} (pōrṭṭar) — a porter}

\label{lesson:ML-C466-porttar}



\begin{quote}
Before the new one: say the Malayalam for a suitcase, then the Malayalam for a load.
\end{quote}

\subsection*{You'll want to know: \ml{പോർട്ടർ}}
\textbf{\ml{പോർട്ടർ}} — \emph{pōrṭṭar} — \textquotedblleft{}a porter\textquotedblright{}.

\begin{quote}
\ml{പോർട്ടർ} \ml{ഞങ്ങളുടെ} \ml{ലഗേജ്} \ml{പ്ലാറ്റ്ഫോമിലേക്ക്} \ml{കൊണ്ടുപോയി}.
\end{quote}

\begin{quote}
\emph{pōrṭṭar ñaṅṅaḷuṭe lagējŭ plāṟṟphōmilēkkŭ koṇṭupōyi.}
\end{quote}

\begin{quote}
The porter took our luggage to the platform.
\end{quote}

At a railway station porters wear a red shirt and a numbered badge.

\begin{cousinweb}
\textbf{\ml{പോർട്ടർ}} (\emph{pōrṭṭar}) was borrowed from English \emph{porter}, a person who carries luggage for travellers.
\end{cousinweb}

\begin{grammarlens}[title={What you've built}]
One more word for getting around.
\end{grammarlens}

\subsection*{Your turn}
\begin{itemize}
\raggedright
  \item \emph{Say it:} \emph{pōrṭṭar}
  \item \emph{Say it:} \emph{pōrṭṭar}, once more
  \item \emph{Say it:} \emph{pōrṭṭar ñaṅṅaḷuṭe lagējŭ plāṟṟphōmilēkkŭ koṇṭupōyi.} — The porter took our luggage to the platform.
  \item \emph{Recall:} say the Malayalam for a suitcase, then the Malayalam for a load, then say \emph{pōrṭṭar} again
\end{itemize}

\begin{tcolorbox}[breakable,colback=teal!4,colframe=teal!35!black,title={Before you move on}]
What does \ml{പോർട്ടർ} mean? (\textquotedblleft{}A porter\textquotedblright{}.) Say it once more.
\end{tcolorbox}
\section[\texorpdfstring{\ml{ക്ലോക്ക്റൂം} (klōkkṟūṁ)}{ക്ലോക്ക്റൂം (klōkkṟūṁ)}]{\ml{ക്ലോക്ക്റൂം} (klōkkṟūṁ) — a cloakroom; a left-luggage room}

\label{lesson:ML-C466-klokkrum}



\begin{quote}
Before the new one: say the Malayalam for a load, then the Malayalam for a porter.
\end{quote}

\subsection*{You'll want to know: \ml{ക്ലോക്ക്റൂം}}
\textbf{\ml{ക്ലോക്ക്റൂം}} — \emph{klōkkṟūṁ} — \textquotedblleft{}a cloakroom; a left-luggage room\textquotedblright{}.

\begin{quote}
\ml{ഞങ്ങൾ} \ml{ബാഗുകൾ} \ml{സ്റ്റേഷനിലെ} \ml{ക്ലോക്ക്റൂമിൽ} \ml{വച്ചു}.
\end{quote}

\begin{quote}
\emph{ñaṅṅaḷ bāgukaḷ sṟṟēṣanile klōkkṟūmil vaccu.}
\end{quote}

\begin{quote}
We left our bags in the station cloakroom.
\end{quote}

You get a receipt, a \ml{രസീത്}, and show it when you take the bag back.

\begin{cousinweb}
\textbf{\ml{ക്ലോക്ക്റൂം}} (\emph{klōkkṟūṁ}) was borrowed from English \emph{cloakroom}, a room where bags and coats are left for a while.
\end{cousinweb}

\begin{grammarlens}[title={What you've built}]
One more word for getting around.
\end{grammarlens}

\subsection*{Your turn}
\begin{itemize}
\raggedright
  \item \emph{Say it:} \emph{klōkkṟūṁ}
  \item \emph{Say it:} \emph{klōkkṟūṁ}, once more
  \item \emph{Say it:} \emph{ñaṅṅaḷ bāgukaḷ sṟṟēṣanile klōkkṟūmil vaccu.} — We left our bags in the station cloakroom.
  \item \emph{Recall:} say the Malayalam for a load, then the Malayalam for a porter, then say \emph{klōkkṟūṁ} again
\end{itemize}

\begin{tcolorbox}[breakable,colback=teal!4,colframe=teal!35!black,title={Before you move on}]
What does \ml{ക്ലോക്ക്റൂം} mean? (\textquotedblleft{}A cloakroom; a left-luggage room\textquotedblright{}.) Say it once more.
\end{tcolorbox}
\section[\texorpdfstring{\ml{വലിച്ചുകൊണ്ടുപോകുക} (valiccukoṇṭupōkuka)}{വലിച്ചുകൊണ്ടുപോകുക (valiccukoṇṭupōkuka)}]{\ml{വലിച്ചുകൊണ്ടുപോകുക} (valiccukoṇṭupōkuka) — to drag along}

\label{lesson:ML-C466-valiccukontupokuka}



\begin{quote}
Before the new one: say the Malayalam for a porter, then the Malayalam for a cloakroom.
\end{quote}

\subsection*{You'll want to know: \ml{വലിച്ചുകൊണ്ടുപോകുക}}
\textbf{\ml{വലിച്ചുകൊണ്ടുപോകുക}} — \emph{valiccukoṇṭupōkuka} — \textquotedblleft{}to drag along\textquotedblright{}.

\begin{quote}
\ml{അവൾ} \ml{ഭാരമുള്ള} \ml{സ്യൂട്ട്കേസ്} \ml{പ്ലാറ്റ്ഫോമിലൂടെ} \ml{വലിച്ചുകൊണ്ടുപോയി}.
\end{quote}

\begin{quote}
\emph{avaḷ bhāramuḷḷa syūṭṭkēsŭ plāṟṟphōmilūṭe valiccukoṇṭupōyi.}
\end{quote}

\begin{quote}
She dragged the heavy suitcase along the platform.
\end{quote}

Past: \ml{വലിച്ചുകൊണ്ടുപോയി}, \emph{valiccukoṇṭupōyi}.

Three verbs in a row: \ml{വലിച്ച്}, having pulled, \ml{കൊണ്ട്}, holding, \ml{പോകുക}, to go.

\begin{cousinweb}
\textbf{\ml{വലിച്ചുകൊണ്ടുപോകുക}} (\emph{valiccukoṇṭupōkuka}) is a native Dravidian compound of \ml{വലിക്കുക}, to pull (DEDR 5282), and \ml{കൊണ്ടുപോകുക}, to take away, built on \ml{കൊള്ളുക}, to hold (DEDR 2151).
\end{cousinweb}

\begin{grammarlens}[title={What you've built}]
One more word for getting around.
\end{grammarlens}

\subsection*{Your turn}
\begin{itemize}
\raggedright
  \item \emph{Say it:} \emph{valiccukoṇṭupōkuka}
  \item \emph{Say it:} \emph{valiccukoṇṭupōkuka}, once more
  \item \emph{Say it:} \emph{avaḷ bhāramuḷḷa syūṭṭkēsŭ plāṟṟphōmilūṭe valiccukoṇṭupōyi.} — She dragged the heavy suitcase along the platform.
  \item \emph{Recall:} say the Malayalam for a porter, then the Malayalam for a cloakroom, then say \emph{valiccukoṇṭupōkuka} again
\end{itemize}

\begin{tcolorbox}[breakable,colback=teal!4,colframe=teal!35!black,title={Before you move on}]
What does \ml{വലിച്ചുകൊണ്ടുപോകുക} mean? (\textquotedblleft{}To drag along\textquotedblright{}.) Say it once more.
\end{tcolorbox}
